\scanpage{137}
\mspage{V.9}
% archive sheet number: 136
% manuscript lines 1-4
\deleted{\unsure{2.4.}} \deleted{Supposons} \inserted{Considérons} maintenant \deleted{que} l'ens \inserted{(fermé) $Z$} des points de $X_0$ où $f:X\mathbin{\unsure{-|\to}}S$ est non lisse, \deleted{\unsure{soit de}} \inserted{$d$ sa} dim. \deleted{\unsure{$d$}}
On a \deleted{alors} \inserted{$\operatorname{supp}\underline\Phi^q_{\unsure{\nu|v}}\subset Z$ (\uline{III} 1.), donc,} par SGA 4 \deleted{\unsure{6}} \uline{X}
\[
H^p(X_0,\underline\Phi^q_\nu)\ne0\quad\deleted{\text{\unsure{signifie}}}\inserted{\text{implique}}\quad p\leq2d.
\tag*{(*)}
\]
% manuscript lines 5-9
D'ailleurs, on a en vertu de 2.\unsure{3}. (où $n=\dim X_{\unsure{\bar\eta|\bar y}}$)
\[
\delta(\underline\Phi^q_\nu)\leq n+1-q\qquad\text{si }q\geq1
\]
donc on trouve une \unsure{valeur|résultat} \inserted{$\delta(\Phi^q_\nu)$} $\leq d$ si $n+1-q<d$ i.e. $q>n+1-d$; \deleted{donc} pour de tels $q$, on a par SGA 4 \uline{X}
\[
\begin{aligned}
&\deleted{2\delta({\unsure{\underline\Phi|\Phi}}^q_\nu)\unsure{.}}\\
&H^p(X_0,\underline\Phi^q_\nu)\ne0\quad\deleted{\text{\unsure{signifie}}}\inserted{\text{implique}}\quad p\leq2(n+1-q)\\
&\hspace{8em}\deleted{\text{i.e. implique}}\qquad p+q\leq2n+2-q,
\end{aligned}
\tag*{(**)}
\]
% manuscript lines 10-14
et on a \inserted{pour ces $q$,} $2n+2-q<2n+2-(n+1-d)=n+1+d$. \unsure{Donc}
D'autre part, si $q\leq n+1-d$ et \inserted{$H^p(X_0,\underline\Phi^q_\nu)\ne0$ i.e.} $p\leq2d$, on a $p+q\leq n+1+d$. Donc en tous cas (si $q\geq1$)
\[
H^p(X_0,\underline\Phi^q_\nu)\ne0\quad\Longrightarrow\quad p+q\leq n+d+1.
\]
% manuscript lines 15-19
\deleted{Supposons} Cette implication est aussi vraie pour $q=0$, si on suppose que $d\leq n+1$ (ce qui est le cas \inserted{(et $d\leq n=\dim X_0$), p. ex. si $X=\overline{X_\eta}$,} \deleted{\unsure{en particulier}} \unsure{i.e.} $X$ plat sur $S$]. On trouve \inserted{alors (\unsure{cf. 1})}
\[
\Phi^i_\nu=0\qquad\text{pour }i>n+d+1.
\]
\inserted{Donc :}

% manuscript lines 20-27
\uline{Corollaire} 2.5. \deleted{Supposons que} Soit $d$ la dimension de l'ens des pts de $X_0$ en lesquels $f:X\mathbin{\unsure{-|\to}}S$ est non lisse, \deleted{et} \inserted{$n=\dim X_{\bar\eta}$, et} supposons que $d\leq n+1$ i.e. que $\dim X_0\leq n+1$. [ce qui est le cas si $X=\overline{X_\eta}$, impliquant $\dim X_0=n$]. Alors le \unsure{grp|grps} $\Phi^i_\nu$ des cycles évanescents \inserted{(\uline{III} 1.)}, à coefficients dans $\Lambda_\nu=\mathbb Z/\ell^\nu\mathbb Z$, est nul pour \deleted{\unsure{$q$}} $i>n+d+1$, en d'autres termes l'hom
% slanted
\margin{! (rel. à $U=X_\eta$)}
